\documentclass[25pt, a0paper, portrait, margin=0mm, innermargin=15mm, blockverticalspace=15mm, colspace=15mm]{tikzposter}

\usepackage{amsmath}
\usepackage{graphicx}
\usepackage{booktabs}

\definecolor{ftteal}{RGB}{13,148,136}
\definecolor{ftindigo}{RGB}{79,70,229}
\definecolorstyle{FreedomColors}{
  \definecolor{colorOne}{named}{ftteal}
  \definecolor{colorTwo}{named}{ftindigo}
  \definecolor{colorThree}{named}{white}
}{
  \colorlet{backgroundcolor}{white}
  \colorlet{framecolor}{colorOne}
  \colorlet{titlefgcolor}{white}
  \colorlet{titlebgcolor}{colorOne}
  \colorlet{blocktitlebgcolor}{colorTwo}
  \colorlet{blocktitlefgcolor}{white}
  \colorlet{blockbodybgcolor}{white}
  \colorlet{blockbodyfgcolor}{black}
  \colorlet{innerblocktitlebgcolor}{colorOne}
  \colorlet{innerblocktitlefgcolor}{white}
  \colorlet{innerblockbodybgcolor}{colorOne!10}
  \colorlet{innerblockbodyfgcolor}{black}
  \colorlet{notefgcolor}{black}
  \colorlet{notebgcolor}{colorOne!20}
  \colorlet{notefrcolor}{colorOne}
}
\usetheme{Simple}
\usecolorstyle{FreedomColors}

\title{Poster Title: A Short and Descriptive Headline}
\author{First Author, Second Author}
\institute{Multimedia University, Malaysia}

\begin{document}

\maketitle

\begin{columns}
  \column{0.5}
  \block{Introduction}{
    Explain the problem and why it matters in two or three sentences. Posters are read from a distance, so keep text short and use large figures.
  }
  \block{Objectives}{
    \begin{itemize}
      \item First objective
      \item Second objective
      \item Third objective
    \end{itemize}
  }
  \block{Method}{
    Describe the approach. Key equation:
    \[
      \hat{y} = \arg\max_{c} \; P(c \mid x)
    \]
    \innerblock{Highlight}{Use inner blocks to emphasise a key idea.}
  }

  \column{0.5}
  \block{Results}{
    \begin{center}
      \begin{tabular}{lcc}
        \toprule
        Method & Accuracy & F1 \\
        \midrule
        Baseline & 0.84 & 0.81 \\
        Proposed & \textbf{0.93} & \textbf{0.90} \\
        \bottomrule
      \end{tabular}
    \end{center}
  }
  \block{Conclusion}{
    Summarise the main take-away message in one or two sentences.
  }
  \note[targetoffsetx=-8cm, targetoffsety=-5cm, width=10cm]{Notes like this one can point at things.}
  \block{References}{
    \small
    [1] A. Author, ``An important paper,'' \textit{Journal}, 2024.
  }
\end{columns}

\end{document}
